% TOP_PROBLEM: {"id": 3122, "title": "The Borodin-Kostochka Conjecture", "category_id": 3, "category": {"id": 3, "name": "graph_theory", "display_name": "Graph Theory", "description": "Problems involving graphs, networks, and their properties.", "slug": "graph-theory", "order_index": 3, "created_at": "2026-07-31T15:26:25.671Z"}, "status": "open", "problem_number": "OPG-37670", "set_id": 12}
% FIRST_POSTED: 2026-10-01
% ATTEMPT_STATUS: unresolved
% UnsolvedMath ID 3122; source catalogue code OPG-37670
% !TeX program = lualatex
% GPT-6 Astra Ultra research attempt, 2026-10-01; self-contained partial work.
% Completion estimate: no numerical percentage assigned; full problem unresolved.
\documentclass[11pt,a4paper]{article}
\usepackage[margin=25mm]{geometry}
\usepackage{fontspec}
\setmainfont{lmroman10-regular.otf}[BoldFont=lmroman10-bold.otf,
 ItalicFont=lmroman10-italic.otf,BoldItalicFont=lmroman10-bolditalic.otf]
\usepackage{amsmath,amssymb,mathtools}
\usepackage{unicode-math}
\setmathfont{latinmodern-math.otf}
\AtBeginDocument{\let\setminus\smallsetminus}
\usepackage{xurl,hyperref}
\hypersetup{colorlinks=true,urlcolor=blue}
\setcounter{secnumdepth}{0}
\setlength{\parindent}{0pt}
\setlength{\parskip}{5pt}
\setlength{\emergencystretch}{3em}
\begin{document}
\section*{The Borodin-Kostochka Conjecture}\label{3122}
{\small UnsolvedMath ID 3122 · OPG-37670 · Graph Theory · OpenGarden}
\subsection{Definitions and mathematical statement}
Let $G=(V,E)$ be a finite simple undirected graph. The degree of a
vertex is its number of neighbours, and $\Delta(G)$ is the largest
degree. A clique is a vertex subset in which every pair is adjacent;
its maximum possible cardinality is $\omega(G)$.
A proper vertex colouring is a map $c:V\to\{1,\ldots,q\}$ assigning
different colours to adjacent vertices, and $\chi(G)$ is the least
such $q$.

The Borodin--Kostochka conjecture states that for every $G$ satisfying
$\Delta(G)\ge9$,
\[
 \chi(G)\le\max\{\Delta(G)-1,\omega(G)\}.
\]
In particular, if $G$ has no clique of size $\Delta(G)$, the conclusion
requires a proper colouring with at most $\Delta(G)-1$ colours. If a
large clique is present, the formula explicitly retains the unavoidable
clique lower bound.

The degree threshold is part of the assertion: no claim is made here
for graphs with maximum degree eight or smaller. The condition is on
maximum degree, not minimum degree or average degree. The graph need
not be connected, regular, planar, or from a particular forbidden-subgraph
class. Ordinary colour palettes are used; requiring a colouring from
prescribed lists at the vertices would be a stronger conjecture.

The target is an existence theorem uniformly over graph size. It does
not merely ask for an estimate that approaches $\Delta(G)$ asymptotically:
the saving of one full colour whenever the clique bound permits it is
the exact conclusion.
\subsection{Short English statement}
Does every graph of maximum degree at least nine need at most the larger of its clique number and one less than its maximum degree colours?
\subsection{Research attempt}
\paragraph{Scope and process.}
GPT-6 Astra Ultra, 1 October 2026. This attempt keeps the catalogue statement above, checks primary literature, and develops a restricted argument with an adversarial gap audit. The proofs below are the mathematical output of this attempt, not a claim of novelty or external certification.

\paragraph{Literature and attack plan.}
Dvořák, Kang and Mikšaník [R1] prove a correspondence-colouring version for maximum-degree bounds at least $3\cdot10^9$, strengthening earlier large-degree results. The threshold in the catalogue is nine, so that result does not by itself settle the statement here. We develop a self-contained reduction for the difficult clique-free case and an exact proof for graphs supplied with a perfect elimination ordering. Kempe-chain obstructions identify the point where a greedy extension ceases to be justified.

\paragraph{Reducing a putative obstruction.}
Fix $D\ge9$ and suppose $\Delta(G)\le D$, $\omega(G)\le D-1$, but $G$ is not $(D-1)$-colourable. Choose an induced subgraph $H$ of minimum order with the latter property and put $q=D-1$. Every $H-v$ has a $q$-colouring. If $d_H(v)\le q-1$, at most $q-1$ colours occur on its neighbours, so a missing colour extends that colouring to $v$, a contradiction. Therefore
\[
 D-1\le\delta(H)\le\Delta(H)\le D. \tag{1}
\]
Moreover, $H$ is connected: otherwise every component, being a proper induced subgraph, could be coloured separately with the same palette.

The reduction also proves a useful subclass result without any structural theorem: if every nonempty induced subgraph of $G$ has a vertex of degree at most $D-2$, repeated removal and reverse greedy colouring give $\chi(G)\le D-1$. Any counterexample to this part of the conjecture must instead contain the high-degree core (1).

\paragraph{The local Kempe constraint.}
Suppose $v\in V(H)$ has degree exactly $q$, and fix a $q$-colouring of $H-v$. Every colour must appear on $N_H(v)$; with exactly $q$ neighbours, each appears exactly once. Write $u_i$ for the neighbour coloured $i$.

For any two different colours $i,j$, the vertices $u_i,u_j$ must belong to the same connected component of the subgraph induced by colours $i,j$. Otherwise interchange $i,j$ in the component containing $u_i$. This preserves properness: edges within that component remain bichromatic, edges to other colours remain proper, and an edge to an $i$- or $j$-coloured vertex outside the component would contradict the definition of component. After swapping, no neighbour of $v$ has colour $i$, because its unique such neighbour changed to $j$ while $u_j$ was not changed. Assigning $i$ to $v$ gives the forbidden $q$-colouring of $H$.

Hence every pair of the uniquely coloured neighbours is joined by a corresponding two-colour path. This is a necessary condition, not a statement that the neighbours form a clique. A two-colour path can have many internal vertices, so replacing ``same component'' by ``adjacent'' would be the first unjustified leap in this strategy.

\paragraph{A complete proof for elimination-ordered graphs.}
Suppose $G$ has an ordering $v_1,\ldots,v_n$ such that the later neighbours of each $v_i$ form a clique. Colour in reverse order. When reaching $v_i$, its already coloured neighbours number at most $\omega(G)-1$, since together with $v_i$ they form a clique. A palette of $\omega(G)$ therefore always has a free colour. This proves
\[
 \chi(G)=\omega(G)
\]
for every graph with such an ordering, and hence proves the catalogue inequality on this subclass for every maximum degree, in particular $\Delta\ge9$. The proof requires the stated ordering; it does not assert that arbitrary graphs admit it. Complete graphs, trees rooted and removed from leaves, and recursively clique-attached examples provide immediate verifiable instances.

\paragraph{Verification.}
The minimality argument does not assume the induced subgraph $H$ still has maximum degree at least nine; it works with the fixed ambient bound $D$ throughout. This avoids incorrectly applying the conjecture recursively to lower-degree subgraphs. The Kempe argument at a degree-$q$ vertex uses uniqueness of each neighbour colour. At a degree-$q+1$ vertex a colour may repeat, so the same swap need not free any colour; that case is not covered. Finally the clique restriction $\omega\le D-1$ was stated explicitly, and no reduction of the entire conjecture to this restricted case is claimed without additional results.

\paragraph{First gap and final status.}
The attempt proves the elimination-order case, the low-degeneracy case, and necessary degree/Kempe constraints on an obstruction in the difficult clique-free regime. It does not exclude cores with degrees $D-1,D$, or convert the two-colour connecting paths into the forbidden clique. The full Borodin--Kostochka conjecture is \textbf{UNRESOLVED} by this attempt.

\subsection{Sources}
\begin{itemize}
\item[S1] ULAM.AI and the UnsolvedMath Contributors, supplied problems.json, record 3122 (OPG-37670), OpenGarden collection. Catalogue identity and source transcription; CC BY 4.0.
\url{https://huggingface.co/datasets/ulamai/UnsolvedMath}.
\item[S2] Open Problem Garden, The Borodin-Kostochka Conjecture. Original problem and discussion; the mathematical formulation below is expanded from the supplied source record.
\url{http://www.openproblemgarden.org/op/the_borodin_kostochka_conjecture}.
\item[R1] Z. Dvořák, R. J. Kang and D. Mikšaník, \emph{On Borodin--Kostochka conjecture for correspondence coloring}, arXiv:2603.14427 (2026). Primary abstract checked for its explicit large-degree threshold.
\url{https://arxiv.org/abs/2603.14427}.
\end{itemize}
\end{document}
